\section{Diskusi}
\label{sec:discussion}

\subsection{Asumsi Kepercayaan dan Keterbatasan}
Kami menyatakan secara eksplisit asumsi kepercayaan berikut agar pengadopsi dapat menalar risiko residual:

\begin{enumerate}
  \item \textbf{Kejujuran administrator} --- \texttt{DEFAULT\_ADMIN\_ROLE} memberikan hak istimewa auditor dan produsen.
  Kompromi kunci admin atau kolusi dengan auditor jahat merusak integritas sertifikasi meskipun log permanen.
  \item \textbf{Ketelitian auditor} --- catatan on-chain membuktikan bahwa alamat auditor menyetujui batch pada waktu tertentu; catatan tidak secara kriptografis membuktikan kepatuhan halal bahan, metode penyembelihan, atau perilaku pemasok.
  HalalChain adalah \emph{jejak audit}, bukan klasifikator halal otomatis.
  \item \textbf{Ketersediaan pinning IPFS} --- dokumen memerlukan pin Pinata (atau self-hosted) yang aktif.
  CID yang tidak dipin atau kedaluwarsa menjadi tidak tersedia off-chain sementara CID on-chain tetap ada, menghasilkan status referensi ``menggantung'' yang harus ditafsirkan konsumen dengan hati-hati.
  \item \textbf{Kelangsungan sequencer L2} --- Base mengandalkan infrastruktur rollup; gangguan berkepanjangan menunda tulisan tetapi tidak mengubah status yang dikomitkan secara diam-diam.
  \item \textbf{Ruang lingkup testnet} --- evaluasi pada Base Sepolia tidak menyiratkan integrasi BPJPH produksi, kesetaraan hukum dengan registrasi SIHALAL, atau kesiapan mainnet tanpa tinjauan tata kelola dan manajemen kunci lebih lanjut.
\end{enumerate}

Keterbatasan ini sebagian membatasi persyaratan desain \emph{Adl} (keadilan) hingga tata kelola auditor multi-tanda tangan atau penugasan peran berbasis DAO diimplementasikan~\cite{gnosis2024safe}.
Versi mendatang dapat mensyaratkan persetujuan auditor \texttt{m}-dari-\texttt{n} untuk verifikasi, menyelaraskan kebijakan on-chain dengan lembaga halal institusional.

\subsection{Perbandingan dengan Alternatif Terpusat}
Basis data terpusat menawarkan kompleksitas operasional lebih rendah, alat SQL yang familiar, dan integrasi alur kerja BPJPH yang mudah~\cite{underwood2016blockchain,viriyasitavat2019blockchain}.
Administrator dapat mencabut akses, memperbaiki kesalahan ketik, dan menjalankan laporan batch tanpa biaya gas.
Namun, model terpusat memungkinkan alterasi catatan retroaktif oleh pengguna berhak dan tidak menyediakan verifikasi pihak ketiga tanpa izin---konsumen harus mempercayai respons API operator.

HalalChain menukar properti ini dengan verifikabilitas publik dan jejak audit tambah-saja.
Produsen dan auditor harus mengelola dompet Ethereum dan ETH (atau ETH yang di-bridge pada Base) untuk tulisan; pinning IPFS menimbulkan biaya berulang atau per-GB.
Bagi produsen UMKM, onboarding bersubsidi (meta-transaksi, paymaster, atau sponsor institusional) mungkin diperlukan untuk adopsi skala besar.

Dibandingkan blockchain enterprise (Hyperledger), HalalChain lebih mengutamakan transparansi daripada privasi transaksi: status batch dapat dibaca dunia.
Sertifikasi halal pada dasarnya adalah \emph{klaim publik}; saluran privat mungkin penting untuk formulasi rahasia dagang tetapi di luar ruang lingkup v1.

\subsection{Implikasi bagi UMKM dan Kebijakan}
Mandat pelabelan halal Indonesia menciptakan permintaan alat ketertelusuran terjangkau.
HalalChain menunjukkan arsitektur referensi yang dapat dipilot oleh lembaga terakreditasi BPJPH, koperasi regional, atau tim KKN universitas tanpa melisensikan platform proprietari.
Integrasi dengan SIHALAL memerlukan akses API resmi, pemetaan hukum status on-chain ke nomor registrasi, dan alur kerja akreditasi auditor---tidak ada yang diimplementasikan pada prototipe.

Pasar ekspor semakin mengharapkan ketertelusuran digital~\cite{marianingsih2024halal}.
Halaman verifikasi publik tertaut QR dapat melengkapi label fisik, meskipun penerimaan regulasi bervariasi menurut negara tujuan.

\subsection{Penelitian Lanjutan}
Arah masa depan meliputi:
\begin{itemize}
  \item verifikasi auditor multi-tanda tangan dan pencabutan terkunci waktu;
  \item pemilihan auditor berbasis DAO dengan modul stake atau reputasi;
  \item deployment Base mainnet dengan dasbor pemantauan biaya;
  \item pengaitan sensor IoT (suhu, GPS) melalui jaringan oracle;
  \item integrasi dengan BPJPH SIHALAL jika API tersedia;
  \item verifikasi formal transisi FSM di \texttt{HalalChain.sol};
  \item studi pengguna dengan produsen UMKM dan konsumen tentang kegunaan onboarding dompet dan verifikasi QR.
\end{itemize}
